% Einheit 2 – Terme zusammenfassen (bank/terme/e2.jsonl, zusammenbau v1.0)
\einheitenkopf[e2][Terme zusammenfassen]{Terme zusammenfassen}
\setcounter{aufgabe}{5}


% test: terme-e2-k4-s11-v2, terme-e2-k5-s7-v2
\begin{kbaufgabe}{Kannst du das schon? Dann weiter zu „Klammern auflösen“}
\begin{kbblock}
\kbteil{a)}{Vereinfache den Term $10b - 5b^2 - 7b$ und berechne seinen Wert für $b = 2$.}{P10 ’25 G}
\item[]\kbraum{3}
\end{kbblock}
\begin{kbblock}
\kbanweisung{Multipliziere.}
\kbteil{b)}{$(-4x) \cdot 3y \cdot (-2) =$ \lbfeld}{}
\end{kbblock}
\end{kbaufgabe}


% erkennung: terme-e2-k1-s0-v1, terme-e2-k1-s0-v2, terme-e2-k1-s0-v3, terme-e2-k1-s0-v4
\begin{kbaufgabe}{Gleichartige Glieder erkennen}
\begin{kbblock}
\kbanweisung{Unterstreiche, was zusammengehört.}
\lbpaar{\kbteil{a)}{$4y$, $7$, $9y$}{}
}{\kbteil{b)}{$6$, $2a$, $8$, $a$}{}}
\end{kbblock}
\begin{kbblock}
\lbpaar{\kbteil{c)}{$5x^2$, $4x$, $3x^2$}{}
}{\kbteil{d)}{$7b$, $2c$, $5b$, $4$}{}}
\end{kbblock}
\end{kbaufgabe}


% erkennung: terme-e2-k2-s0-v1, terme-e2-k2-s0-v2, terme-e2-k2-s0-v3, terme-e2-k2-s0-v4
\begin{kbaufgabe}{Die Vorzahl eines Gliedes ablesen}
\begin{kbblock}
\lbpaar{\kbteil{a)}{Welche Vorzahl hat $a$?}{}
\kbantwort{\kbfeld}
}{\kbteil{b)}{Welche Vorzahl hat $-b$?}{}
\kbantwort{\kbfeld}}
\end{kbblock}
\begin{kbblock}
\lbpaar{\kbteil{c)}{Welche Vorzahl hat $1{,}5z$?}{}
\kbantwort{\kbfeld}
}{\kbteil{d)}{Welche Vorzahl hat $-4w$?}{}
\kbantwort{\kbfeld}}
\end{kbblock}
\end{kbaufgabe}


% erkennung: terme-e2-k3-s0-v1, terme-e2-k3-s0-v2, terme-e2-k3-s0-v3, terme-e2-k3-s0-v4
\begin{kbaufgabe}{Die Glieder eines Terms mit ihrem Vorzeichen aufschreiben}
\begin{kbblock}
\kbteil{a)}{Schreibe die Glieder des Terms $5a - 2b + 7$ mit ihrem Vorzeichen auf.}{}
\kbantwort{\kbfeld}
\end{kbblock}
\begin{kbblock}
\kbteil{b)}{Schreibe die Glieder des Terms $-3x + 6 - y$ mit ihrem Vorzeichen auf.}{}
\kbantwort{\kbfeld}
\end{kbblock}
\begin{kbblock}
\kbteil{c)}{Schreibe die Glieder des Terms $2m - 9 - 5k$ mit ihrem Vorzeichen auf.}{}
\kbantwort{\kbfeld}
\end{kbblock}
\begin{kbblock}
\kbteil{d)}{Schreibe die Glieder des Terms $-6 + 3c - 4d$ mit ihrem Vorzeichen auf.}{}
\kbantwort{\kbfeld}
\end{kbblock}
\end{kbaufgabe}


% leiter: terme-e2-k4-s1-v1, terme-e2-k4-s1-v2, terme-e2-k4-s1-v3, terme-e2-k4-s1-v4, terme-e2-k4-s2-v1, terme-e2-k4-s3-v1, terme-e2-k4-s4-v1, terme-e2-k4-s5-v1, terme-e2-k4-s6-v1, terme-e2-k4-s7-v1, terme-e2-k4-s8-v1, terme-e2-k4-s9-v1, terme-e2-k4-s10-v1
\begin{kbaufgabe}{Terme zusammenfassen}
\begin{kbblock}
\kbanweisung{Fasse zusammen.}
\lbpaar{\kbteil{a)}{$7x + 2x =$ \lbfeld}{}
}{\kbteil{b)}{$7x + 4x =$ \lbfeld}{}}
\end{kbblock}
\begin{kbblock}
\lbpaar{\kbteil{c)}{$7x + 5x =$ \lbfeld}{}
}{\kbteil{d)}{$7x + 8x =$ \lbfeld}{}}
\end{kbblock}
\begin{kbblock}
\lbpaar{\kbteil{e)}{$9x - 4x =$ \lbfeld}{}
}{\kbteil{f)}{$2x + 4x + 3x =$ \lbfeld}{}}
\end{kbblock}
\begin{kbblock}
\lbpaar{\kbteil{g)}{$5x + 3 + 2x =$ \lbfeld}{}
}{\kbteil{h)}{$4a + 3b + 2a =$ \lbfeld}{}}
\end{kbblock}
\begin{kbblock}
\lbpaar{\kbteil{i)}{$y + 4y =$ \lbfeld}{}
}{\kbteil{j)}{$2x - 7x =$ \lbfeld}{}}
\end{kbblock}
\begin{kbblock}
\lbpaar{\kbteil{k)}{$-3x + 8x =$ \lbfeld}{}
}{\kbteil{l)}{$4x^2 + 5x + 2x^2 =$ \lbfeld}{}}
\end{kbblock}
\begin{kbblock}
\lbpaar{\kbteil{m)}{$1{,}5x + 3{,}5x =$ \lbfeld}{}
}{\item[]}
\end{kbblock}
\end{kbaufgabe}


% pruefung: terme-e2-k4-s11-v1
\begin{kbaufgabe}{Auch Aufgaben aus der Prüfung zum Zusammenfassen lösen}
\begin{kbblock}
\kbteil{}{Vereinfache den Term $8a - 2a^2 - 5a$ und berechne seinen Wert für $a = 3$.}{P10 ’25 G}
\item[]\kbraum{3}
\end{kbblock}
\end{kbaufgabe}


% leiter: terme-e2-k5-s1-v1, terme-e2-k5-s1-v2, terme-e2-k5-s1-v3, terme-e2-k5-s1-v4, terme-e2-k5-s2-v1, terme-e2-k5-s3-v1, terme-e2-k5-s4-v1, terme-e2-k5-s5-v1, terme-e2-k5-s6-v1
\begin{kbaufgabe}{Terme malnehmen}
\begin{kbblock}
\kbanweisung{Multipliziere.}
\lbpaar{\kbteil{a)}{$5 \cdot 2x =$ \lbfeld}{}
}{\kbteil{b)}{$5 \cdot 3x =$ \lbfeld}{}}
\end{kbblock}
\begin{kbblock}
\lbpaar{\kbteil{c)}{$5 \cdot 6x =$ \lbfeld}{}
}{\kbteil{d)}{$5 \cdot 7x =$ \lbfeld}{}}
\end{kbblock}
\begin{kbblock}
\lbpaar{\kbteil{e)}{$4x \cdot 5 =$ \lbfeld}{}
}{\kbteil{f)}{$4y \cdot 2y =$ \lbfeld}{}}
\end{kbblock}
\begin{kbblock}
\lbpaar{\kbteil{g)}{$(-4) \cdot 5x =$ \lbfeld}{}
}{\kbteil{h)}{$4a \cdot 5b =$ \lbfeld}{}}
\end{kbblock}
\begin{kbblock}
\lbpaar{\kbteil{i)}{$2 \cdot 5x \cdot 3 =$ \lbfeld}{}
}{\item[]}
\end{kbblock}
\end{kbaufgabe}


% pruefung: terme-e2-k5-s7-v1
\begin{kbaufgabe}{Auch Produkte mit Vorzeichen und zwei Variablen berechnen}
\begin{kbblock}
\kbanweisung{Multipliziere.}
\kbteil{}{$(-6a) \cdot 2b =$ \lbfeld}{}
\end{kbblock}
\end{kbaufgabe}


% pflicht: terme-e2-k6-s1-v1, terme-e2-k6-s2-v1, terme-e2-k6-s3-v1, terme-e2-k6-s4-v1
\begin{kbaufgabe}{Verstanden?}
\begin{kbblock}
\kbteil{a)}{Mia soll den Term $5a + 3 + 2a$ zusammenfassen. Sie rechnet so: \rechnung{5a + 3 + 2a &= 10a} Finde den Fehler und rechne richtig.}{}
\item[]\kbraum{2}
\end{kbblock}
\begin{kbblock}
\kbteil{b)}{Begründe, ob man den Term $4x + 4y$ zusammenfassen kann.}{}
\item[]\kbraum{3}
\end{kbblock}
\begin{kbblock}
\kbteil{c)}{Die Figur zeigt ein Rechteck. Die Maße sind in cm angegeben. Schreibe einen Term für den Umfang. Fasse den Term zusammen.}{}
\kbzweispaltig[0.48]{\viereck[seiten={3x,2,3x,2}]{(0,0)}{(6,0)}{(6,2)}{(0,2)}}{\lbantwortrechts[1.2cm]{\kbfeld}}
\end{kbblock}
\begin{kbblock}
\kbteil{d)}{Ein Kiosk verkauft Brötchen zu je $p$ Euro. Am Montag verkauft er 45 Stück, am Dienstag 38 und am Mittwoch 52. Stelle einen Term für die Einnahmen der drei Tage auf und fasse ihn zusammen. Der Kiosk will in den drei Tagen mindestens 50 € mit Brötchen einnehmen. Schafft er das bei $p = 0{,}40$?}{}
\item[]\kbraum{3}
\end{kbblock}
\end{kbaufgabe}
